\documentclass[instructions.tex]{subfiles}
\begin{document}
\chapter{Le tort que l’homme se fait à lui-même par le péché mortel.}


Tout homme qui pèche mortellement, est un aveugle\footnote{Omnis peccans est ignorans.}. Il ne comprend ni ce qu’il fait, ni ce qu’il perd. Il perd la grâce sanctifiante, il perd son âme, il perd son Dieu.


\primo En perdant la grâce sanctifiante, il perd un don mille fois plus précieux que tous les trésors de la terre. Avec elle l’homme le plus pauvre est heureux, et sans elle l’homme le plus riche est misérable; parce que c’est la grâce sanctifiante qui rend l’homme ami de Dieu et digne du ciel : l’homme dépouillé de la grâce sanctifiante est donc ennemi de Dieu et exclu de son royaume. Être ennemi de Dieu ! Cette pensée fait frémir. Oh ! si l’on craint d’être disgracié d’un roi de la terre, comment peut-on être tranquille en vivant dans la disgrâce d’un Dieu ?

\secundo La perte de la grâce entraîne la perte des mérites qu’on avait acquis. Plusieurs années de combats et de victoires dans l’exercice de la patience et des autres vertus vous avaient coûté des peines, et vous avaient acquis des mérites pour le ciel. Si vous êtes assez lâche pour tomber dans un péché mortel, toutes ces peines et tous ces mérites sont perdus : perte plus grande, et que vous devez plus déplorer que si vous perdiez tous les royaumes de l’univers. Vous pleurez la perte d’un argent périssable, dit saint Augustin ; pleurez bien plutôt la perte de vos mérites, puisqu’en les perdant vous perdez tout.

\tertio Étant privé de la grâce, on n’est pas seulement privé des mérites qu’on a acquis, on est encore hors d’état d’en acquérir. La consolation d’un fidèle dans l’état de grâce est que toutes ses actions, quand il les rapporte à Dieu, lui acquièrent de nouveaux mérites. Mais, quand vous feriez les bonnes œuvres de tous les saints, si elles sont faites en péché mortel, elles sont perdues pour le ciel.

N’en concluez pas qu’il est inutile de faire de bonnes œuvres dans ce malheureux état ; loin de là, vous devez en faire, et en faire beaucoup. Quoique ces bonnes œuvres, quand on les fait en péché mortel, soient sans valeur pour le ciel, elles sont cependant commandées ; elles sont utiles, et sont même nécessaires pour vous disposer à vous convertir et à vous réconcilier avec Dieu. Tâchez donc de rentrer dans sa grâce par la pénitence ; et étant réconcilié avec lui, vous serez alors en état d’acquérir des mérites pour le ciel.


\section{Résolutions}

\primo Je considérerai donc encore dans le péché mortel les torts immenses qu’il m’a faits à moi-même : il m’a ravi la grâce qui me rendait ami de Dieu, enlevé les mérites que j’avais acquis pour le ciel, m’a réduit dans la triste impossibilité de rien faire qui fût digne d’être récompensé dans l’autre vie, pendant que j’étais dans l’état du péché mortel.

Mon Dieu, que je suis devenu pauvre et hideux devant vous ! Ah ! lavez-moi avec le sang de mon Sauveur, faites que je recouvre les mérites perdus, et que j’en acquière de nouveaux.
